%% =====================================================================
%%  Capítulo 1 -- Introducción y motivación
%%  Formato nuevo (noviembre de 2026). Integra: Handout1.tex (2022), la
%%  reescritura Handout1_2026 (agosto de 2026, con ideas de Kellett y
%%  Braun: modelos de estado, fenómenos sin contraparte lineal, ejemplos
%%  multidominio, integrador de Brockett, conclusiones plausibles y falsas)
%%  y las diapositivas Intro2026.pptx.
%%  Correcciones respecto a las versiones anteriores: ver \S\ref{sec:in-cierre}.
%%  Números verificados con herramientas/cap01_numeros.py (36/36) y
%%  codigo/cap01_introduccion/verificarCap01.m (34/34 en Octave 8.4).
%% =====================================================================

\capitulo[cap:introduccion]{Introducción y motivación}
  {No solo en la investigación, sino también en el mundo cotidiano de la política y la economía, todos estaríamos mejor si más gente se diera cuenta de que los sistemas no lineales simples no poseen necesariamente propiedades dinámicas simples.}
  {Robert May}

\begin{objetivos}[partes/01-introduccion.html]
  \item Escribir un modelo en variables de estado $\dot x=f(t,x,u)$, $y=h(t,x,u)$, encontrar sus equilibrios y trasladarlos al origen.
  \item Reconocer cinco comportamientos que ningún sistema lineal puede tener: equilibrios aislados múltiples, ciclos límite, escape en tiempo finito, dependencia de la amplitud y caos.
  \item Explicar con un ejemplo por qué el todo no es la suma de las partes, y por qué la versión discreta y la continua de la ecuación logística no se comportan igual.
  \item Usar la energía para concluir algo sobre un sistema sin resolverlo, y decir qué \emph{no} se puede concluir todavía.
  \item Dar ejemplos en los que el control lineal no basta, incluido uno en el que ningún control continuo funciona (Brockett).
  \item Distinguir una conclusión plausible de una correcta, que es el hilo de todo el curso.
\end{objetivos}

%% ---------------------------------------------------------------------
\section{¿Por qué sistemas no lineales?}
\label{sec:in-porque}
%% ---------------------------------------------------------------------

Un \clave{sistema dinámico} es una regla que dice cómo cambia en el tiempo un conjunto de variables, las \clave{variables de estado}, que resumen todo lo que hace falta saber del pasado para predecir el futuro. Un modelo dinámico tiene tres ingredientes: un \clave{espacio de estado} donde viven esas variables; una \clave{dinámica}, la regla de evolución; y una \clave{condición inicial} desde la que se arranca. Con esos tres ingredientes se modela un péndulo, un circuito, un tanque, una población de conejos o una economía.

Casi todo lo que usted ha estudiado de sistemas dinámicos es lineal. Un sistema es lineal si cumple el \clave{principio de superposición}: la respuesta a una suma de causas es la suma de las respuestas a cada causa. Esa propiedad es tan poderosa que organiza toda la teoría que usted conoce: la transformada de Laplace, la función de transferencia, el diagrama de Bode, la respuesta al escalón, la ubicación de polos. Un sistema lineal se entiende entendiendo sus partes y sumando. Por eso se dice que la linealidad es el sueño del reduccionista \citep{Mit09}.

\notamargen{\textbf{Ruta de las clases.} Clase~1: \S\ref{sec:in-porque}--\ref{sec:in-poblaciones}. La \S\ref{sec:in-energia}--\ref{sec:in-trabajo} se lee por cuenta propia antes de la clase~2 (capítulo~\ref{cap:unidimensionales}).}

El mundo, en cambio, casi nunca es lineal. Un resorte se endurece al estirarse, un actuador se satura, un diodo conduce en un solo sentido, una población no crece sin límite, la fuerza de un electroimán va como el cuadrado de la corriente. Los modelos lineales funcionan muy bien \emph{cerca} de un punto de operación; el curso trata de lo que pasa lejos de él, y de lo que se pierde al ignorarlo.

\begin{ejemplo}[label=ej:in-islas]{Veinte conejos y dos islas}
Una población de conejos se reproduce una vez al año. En un primer modelo, cada pareja tiene cuatro crías y muere: la población se duplica, $P[k+1]=2P[k]$. En un segundo modelo, la sobrepoblación mata a algunas crías antes de que se reproduzcan; con una capacidad de $K=32$ conejos y un factor de crecimiento $R=1.6$,
\begin{equation}
  P[k+1]=F(P[k]),\qquad F(P)=R\,P\Bigl(1-\frac{P}{K}\Bigr).
  \label{eq:in-logistica-islas}
\end{equation}
¿Cuántos conejos hay al año siguiente si se ponen 20 en una isla? ¿Y si se reparten 10 y 10 en dos islas iguales?
\solucion Con el modelo lineal da igual: $2\cdot20=40=2\cdot10+2\cdot10$. Con el logístico, $F(20)=1.6\cdot20\cdot\tfrac{12}{32}=12$, mientras que $F(10)+F(10)=2\cdot1.6\cdot10\cdot\tfrac{22}{32}=22$. Repartir la población cambia el resultado (\cref{fig:in-superposicion}): el todo no es la suma de las partes. Además, el mejor reparto es 10 y 10: $F(a)+F(20-a)$ es una parábola en $a$ con máximo en $a=10$.
\end{ejemplo}

\begin{figure}[ht]
\centering
\includegraphics[width=\linewidth]{cap01/f_superposicion}
\caption{Izquierda: el mapa de un año, lineal y logístico. Derecha: conejos al año siguiente según cómo se repartan 20 entre dos islas. En el modelo lineal el reparto no importa.}
\label{fig:in-superposicion}
\end{figure}

\codigorepo{cap01_introduccion/C1_Ej1_Superposicion.m}{islas}{Una isla o dos: superposición}

\begin{ideaclave}
Un sistema no lineal no se entiende por partes. Las herramientas del curso reemplazan la superposición por \emph{geometría}: en lugar de resolver la ecuación, se estudia la forma del campo vectorial, dónde están los equilibrios, qué regiones atrapan a las soluciones y qué funciones (como la energía) decrecen a lo largo de ellas.
\end{ideaclave}

Este capítulo responde cinco preguntas: ¿cómo se escribe un modelo no lineal y qué es un equilibrio (\cref{sec:in-modelos})? ¿Qué puede hacer un sistema no lineal que uno lineal no puede (\cref{sec:in-fenomenos})? ¿Por qué se puede concluir algo sin resolver la ecuación (\cref{sec:in-energia})? ¿Qué aportan los sistemas no lineales al control (\cref{sec:in-control})? ¿Cómo se llegó a las herramientas que vamos a usar, y cómo se trabaja con ellas sin caer en conclusiones plausibles pero falsas (\cref{sec:in-historia,sec:in-trabajo})?

\begin{pausa}
¿Es lineal el sistema $\dot x=-x+1$?
\respuesta{No cumple la superposición: si $x_a$ y $x_b$ son soluciones, $\frac{\dif}{\dif t}(x_a+x_b)=-(x_a+x_b)+2$, que no es la ecuación. Es \emph{afín}. Pero con $z=x-1$ queda $\dot z=-z$, que sí es lineal. La no linealidad que importa es la que no se quita con un cambio de coordenadas: trasladar un equilibrio al origen (\cref{sec:in-modelos}) es gratis.}
\end{pausa}

%% ---------------------------------------------------------------------
\section{Modelos de estado y equilibrios}
\label{sec:in-modelos}
%% ---------------------------------------------------------------------

En todo el curso un sistema no lineal se escribe en \clave{variables de estado}:
\begin{equation}
  \dot x=f(t,x,u),\qquad x(t_0)=x_0,\qquad y=h(t,x,u),
  \label{eq:in-general}
\end{equation}
con $x\in\D\subseteq\R^n$ el estado, $u\in\R^m$ la entrada, $y\in\R^p$ la salida y $\D$ un dominio (abierto y conexo). Componente a componente, $\dot x_i=f_i(t,x_1,\dots,x_n,u_1,\dots,u_m)$, $i=1,\dots,n$.
\notamargen{\textbf{Notación} (la de \citet{Kha02}). El tiempo es el \emph{primer} argumento; los vectores no van en negrilla; $\eq$ es un equilibrio. Las convenciones completas están en ``Cómo leer estas notas''.}
El caso lineal es $\dot x=A(t)x+B(t)u$, $y=C(t)x+D(t)u$.

\begin{definicion}[label=def:in-invariante]{Invariante y variante en el tiempo; forzado}
El sistema \eqref{eq:in-general} es \clave{invariante en el tiempo} si $f$ y $h$ no dependen explícitamente de $t$, y \clave{variante en el tiempo} si dependen. Es \clave{forzado} si tiene entrada y \clave{no forzado} si no la tiene (o si $u\equiv0$).
\end{definicion}

\notamargen{En la literatura clásica se dice \emph{autónomo} y \emph{no autónomo}. En estas notas no: ``autónomo'' se confunde con ``sin entrada'', y lo que importa es si $f$ depende de $t$.}
La entrada admite dos lecturas. Si $u=u(t)$ es una señal conocida (una referencia, una perturbación), el sistema se vuelve no forzado y variante en el tiempo, $\dot x=f(t,x,u(t))$. Si $u=k(t,x)$ es una \clave{ley de control} que depende del estado, el \clave{lazo cerrado}
\begin{equation}
  \dot x=f\bigl(t,x,k(t,x)\bigr)\eqqcolon f_{\mathrm{lc}}(t,x)
  \label{eq:in-lazo}
\end{equation}
tampoco tiene entrada, pero ahora el campo vectorial lo \emph{diseñamos} nosotros. Esa es, en una línea, la diferencia entre \clave{análisis} (estudiar $\dot x=f(t,x)$) y \clave{síntesis} (escoger $k$ para que el lazo cerrado tenga las propiedades deseadas), y es el hilo que une la primera mitad del curso con la segunda. Por eso casi toda la teoría se enuncia para $\dot x=f(t,x)$ o $\dot x=f(x)$.

Una ecuación de orden $n$ se vuelve un sistema de primer orden con $x_1=y$, $x_2=\dot y$, \dots, $x_n=y^{(n-1)}$: si $y^{(n)}=\varphi(y,\dot y,\dots,y^{(n-1)},u)$, entonces
\[
  \dot x_1=x_2,\quad\dots,\quad\dot x_{n-1}=x_n,\quad\dot x_n=\varphi(x_1,\dots,x_n,u).
\]

\begin{definicion}[label=def:in-equilibrio]{Punto de equilibrio}
Un punto $\eq\in\D$ es un \clave{equilibrio} de $\dot x=f(x)$ si $f(\eq)=0$; de $\dot x=f(t,x)$, si $f(t,\eq)=0$ para todo $t\ge t_0$. Para $\dot x=f(x,u)$, el par $(\eq,u^*)$ es un \clave{par de equilibrio} si $f(\eq,u^*)=0$. Un equilibrio es \clave{aislado} si tiene una vecindad sin otros equilibrios.
\end{definicion}

Equivalentemente, $\eq$ es un equilibrio si la función constante $x(t)\equiv\eq$ es una solución. Para $\dot x=Ax$ los equilibrios forman el espacio nulo de $A$: si $A$ es invertible el origen es el único, y si no, hay un subespacio entero de equilibrios, ninguno aislado. \alerta{Un sistema lineal no puede tener dos equilibrios aislados.}

\paragraph{Trasladar el equilibrio al origen.} Si $\eq$ es un equilibrio de $\dot x=f(x)$, con $z=x-\eq$,
\[
  \dot z=f(z+\eq)\eqqcolon g(z),\qquad g(0)=f(\eq)=0 .
\]
Lo mismo con un par $(\eq,u^*)$ y $v=u-u^*$. Por eso los teoremas se enuncian para el origen: la traslación no cambia la forma de las trayectorias, solo su ubicación.
\notamargen{Gratis para las matemáticas, no para la interpretación: al final del diseño hay que deshacer la traslación para saber qué significa ``$z=0$'' en las unidades de la planta. Olvidarlo es un error frecuente al reportar simulaciones.}

\begin{trampa}[label=tr:in-tiempo]{Convertir el tiempo en un estado}
\plausible{Todo sistema variante en el tiempo $\dot x=f(t,x)$ es en realidad invariante: basta agregar el estado $x_{n+1}=t$, con $\dot x_{n+1}=1$.}
Como reescritura es correcta, pero no sirve para estudiar estabilidad: el sistema aumentado $\bigl(\dot x,\dot x_{n+1}\bigr)=\bigl(f(x_{n+1},x),1\bigr)$ \emph{no tiene equilibrios}, porque $\dot x_{n+1}=1\neq0$, y todas sus soluciones son no acotadas en la última coordenada. La pregunta ``¿es estable el origen de $\dot x=f(t,x)$?'' no tiene traducción en el sistema aumentado.
\correcto{Los sistemas variantes en el tiempo necesitan una teoría propia, con nociones de \emph{uniformidad} respecto al instante inicial (capítulo~9).}
\end{trampa}

\paragraph{Una advertencia que atraviesa el curso.} Todo lo anterior supone que \eqref{eq:in-general} tiene sentido: que por cada condición inicial pasa una solución, una sola, y que existe mientras la necesitamos. No siempre es así (\cref{sec:in-fenomenos} y capítulo~\ref{cap:unidimensionales}); las condiciones que lo garantizan (Lipschitz) son el tema del \cref{cap:preliminares}. Por ahora la moraleja es: \alerta{antes de aplicar una herramienta a un modelo, verifique sus hipótesis.}

%% ---------------------------------------------------------------------
\subsection{Reescalar: menos parámetros, más entendimiento}
\label{sec:in-reescalar}
%% ---------------------------------------------------------------------

Así como se traslada un equilibrio, se puede \clave{reescalar} el tiempo y las variables para eliminar parámetros \citep{KeBr23}. Considere el masa-resorte con fricción, $m\ddot y+c\dot y+ky=0$, con $x_1=y$, $x_2=\dot y$. Con un tiempo nuevo $\tau=\omega t$ y $z_1=x_1$, $z_2=\dif z_1/\dif\tau=x_2/\omega$,
\[
  \frac{\dif z_1}{\dif\tau}=z_2,\qquad \frac{\dif z_2}{\dif\tau}=-\frac{k}{m\omega^2}\,z_1-\frac{c}{m\omega}\,z_2 .
\]
Como $\omega$ es libre, se escoge $\omega=\sqrt{k/m}$, y con $\alpha=c/\sqrt{km}$ queda
\begin{equation}
  \frac{\dif z_1}{\dif\tau}=z_2,\qquad\frac{\dif z_2}{\dif\tau}=-z_1-\alpha z_2 .
  \label{eq:in-normalizado}
\end{equation}
Tres parámetros $(m,k,c)$ se reducen a uno, $\alpha$: un barrido que era tridimensional ahora es una recta (\cref{fig:in-reescalamiento}). Con el péndulo $m\ell\ddot\theta=-mg\sin\theta-k\ell\dot\theta$, el mismo procedimiento con $\omega=\sqrt{g/\ell}$ y $\alpha=k/(m\omega)$ da
\begin{equation}
  \frac{\dif z_1}{\dif\tau}=z_2,\qquad\frac{\dif z_2}{\dif\tau}=-\sin z_1-\alpha z_2 .
  \label{eq:in-pendulo-norm}
\end{equation}
Cerca de $z_1=0$, $\sin z_1\approx z_1$ y \eqref{eq:in-pendulo-norm} se parece a \eqref{eq:in-normalizado}: \emph{cerca del equilibrio de abajo, el péndulo es un masa-resorte}. Esa frase, hecha precisa, es el método indirecto de Lyapunov (capítulo~7).

\begin{figure}[ht]
\centering
\includegraphics[width=\linewidth]{cap01/f_reescalamiento}
\caption{Tres masa-resorte con $(m,k,c)$ distintos y el mismo $\alpha=c/\sqrt{km}=0.5$. En la variable $\tau=\omega t$ las tres respuestas son la misma curva.}
\label{fig:in-reescalamiento}
\end{figure}

\codigorepo{cap01_introduccion/C1_Ej5_Reescalamiento.m}{reescalar}{Tres sistemas, una curva}

\begin{cuidado}[Deshacer el reescalamiento]
Una ganancia calculada en coordenadas normalizadas no es la que se programa en el equipo, y un tiempo de establecimiento en $\tau$ no está en segundos. Reportar resultados adimensionales sin la transformación inversa es un error recurrente en los proyectos.
\end{cuidado}

%% ---------------------------------------------------------------------
\section{Cinco cosas que un sistema lineal no puede hacer}
\label{sec:in-fenomenos}
%% ---------------------------------------------------------------------

Un sistema lineal $\dot x=Ax$ es, en cierto sentido, aburrido. O el origen es su único equilibrio, o tiene un subespacio entero de ellos. Sus soluciones son combinaciones de exponenciales y senoides, $x(t)=e^{At}x_0$, que existen para todo $t$. Si oscila, la amplitud la fija la condición inicial y el período no depende de ella. Y si $x_a$ y $x_b$ son soluciones, $\alpha x_a+\beta x_b$ también lo es. Un sistema no lineal puede romper cada una de esas reglas. La lista que sigue es el mapa del curso (\cref{fig:in-panorama}).

\begin{figure*}[t]
\centering
\includegraphics[width=\linewidth]{cap01/f_panorama}
\caption{Cinco comportamientos imposibles para $\dot x=Ax$. (a) Péndulo con fricción: infinitos equilibrios, unos estables y otros no. (b) Van der Pol: el mismo ciclo desde adentro (azul) y desde afuera (violeta). (c) Escape en tiempo finito. (d) Silla-nodo: dos equilibrios chocan y desaparecen. (e)--(f) Lorenz: dos condiciones iniciales a $10^{-9}$ se separan.}
\label{fig:in-panorama}
\end{figure*}

\subsection{Equilibrios aislados múltiples}

El péndulo con fricción, con $x_1=\theta$ medido desde la vertical hacia abajo y $x_2=\dot\theta$,
\begin{equation}
  \dot x_1=x_2,\qquad \dot x_2=-\frac g\ell\sin x_1-\frac km\,x_2 ,
  \label{eq:in-pendulo}
\end{equation}
tiene equilibrios en $\eq=(n\pi,0)$, $n\in\Z$. Infinitos, y de dos clases: con $n$ par el péndulo cuelga, y vuelve si se le perturba; con $n$ impar está invertido, y se cae. La consecuencia es inmediata: \alerta{en un sistema no lineal la estabilidad es una propiedad de un equilibrio, no del sistema}. Preguntar ``¿es estable esta planta?'' es una pregunta mal planteada. Y aparece una pregunta que en lo lineal no existe: no basta saber a dónde converge la solución, hay que saber \emph{desde dónde}: la \clave{región de atracción} (capítulos~\ref{cap:lyapunov} y~8).

\subsection{Ciclos límite: oscilar sin que nadie lo sostenga}

El oscilador armónico $\ddot x+x=0$ oscila, pero la amplitud la decide la condición inicial: hay una familia continua de órbitas cerradas, y cualquier perturbación lo pasa a otra. Compare con el oscilador de van der Pol \citep{van26},
\begin{equation}
  \ddot x-\varepsilon\,(1-x^2)\,\dot x+x=0,\qquad\varepsilon>0 .
  \label{eq:in-vdp}
\end{equation}
El término $-\varepsilon(1-x^2)\dot x$ es un amortiguamiento que \emph{cambia de signo}: inyecta energía cuando $\abs x<1$ y la disipa cuando $\abs x>1$. El resultado es una órbita cerrada \emph{aislada}, un \clave{ciclo límite}, a la que llegan las soluciones desde adentro y desde afuera; con $\varepsilon=1$, su amplitud es $2.009$ desde cualquier condición inicial distinta del origen (\cref{fig:in-panorama}b). La amplitud y la frecuencia son propiedades \emph{del sistema}. El latido del corazón, los osciladores de radio, los relojes mecánicos y las neuronas son ciclos límite; si fueran osciladores lineales, un empujón les cambiaría la amplitud para siempre. Detectarlos, y descartarlos, es el tema del \cref{cap:orbitas}.

\codigorepo{cap01_introduccion/C1_Ej3_Panorama.m}{vanderpol}{El mismo ciclo desde adentro y desde afuera}

\subsection{Escape en tiempo finito}

Las soluciones de un sistema lineal inestable crecen, pero tardan un tiempo infinito en llegar a infinito. Un sistema no lineal puede llegar \emph{a una hora fija}. Con $\dot x=-x^2$, $x(0)=-1$, separando variables, $x(t)=1/(t-1)$: la solución tiende a $-\infty$ cuando $t\to1^-$ y \emph{no existe} después (\cref{fig:in-panorama}c). Un integrador numérico no lo anuncia; reporta números cada vez más grandes. Un pariente todavía más incómodo es la \clave{no unicidad}: $\dot x=x^{1/3}$ con $x(0)=0$ tiene infinitas soluciones. Los dos fenómenos se ven en el \cref{cap:unidimensionales} y se explican en el \cref{cap:preliminares}.

\subsection{La amplitud importa: período, armónicos, sincronización}

En un sistema lineal invariante en el tiempo, una entrada senoidal produce en régimen permanente una salida senoidal de la \emph{misma} frecuencia; sobre ese hecho descansa el diagrama de Bode. En un sistema no lineal la salida tiene armónicos, a veces subarmónicos, y el período de una oscilación libre depende de su amplitud. El péndulo sin fricción con amplitud $\theta_0$ tiene período
\[
  T(\theta_0)=\frac{4}{\omega_0}\,K\!\left(\sin\frac{\theta_0}{2}\right),\qquad K(k)=\int_0^{\pi/2}\frac{\dif\varphi}{\sqrt{1-k^2\sin^2\varphi}},
\]
donde $\omega_0=\sqrt{g/\ell}$ y $T_0=2\pi/\omega_0$ es el período de las oscilaciones pequeñas. La \cref{tab:in-periodo} muestra que el ``período del péndulo'' solo existe para amplitudes pequeñas. Y hay fenómenos sin ninguna contraparte lineal, como la \clave{sincronización}: una población de osciladores acoplados que, por encima de un acople crítico, adopta espontáneamente una frecuencia común \citep{Kur75}. Es el modelo de Kuramoto del \cref{cap:orbitas}, el mismo que aparece, con otras unidades, en el sincronismo de los generadores de un sistema de potencia.

\begin{table}[ht]
\centering
\begin{tabular}{@{}lcccc@{}}
\toprule
Amplitud $\theta_0$ & $10\grados$ & $45\grados$ & $90\grados$ & $170\grados$\\
\midrule
$T/T_0$ & $1.0019$ & $1.0400$ & $1.1803$ & $2.4394$\\
\bottomrule
\end{tabular}
\caption{Período del péndulo sin fricción contra la amplitud (\cmd{C1\_Ej3\_Panorama.m}).}
\label{tab:in-periodo}
\end{table}

\codigorepo{cap01_introduccion/C1_Ej3_Panorama.m}{periodo}{El período depende de la amplitud}

\subsection{Bifurcaciones y caos}

Los cuatro fenómenos anteriores ocurren con parámetros fijos. El quinto tiene que ver con lo que pasa cuando los parámetros \emph{cambian}. En $\dot x=r+x^2$, para $r<0$ hay dos equilibrios, $\pm\sqrt{-r}$, uno estable y otro inestable; en $r=0$ chocan; para $r>0$ no hay ninguno y todas las soluciones escapan. El retrato de fase no se deforma: se \emph{reorganiza}. Eso es una \clave{bifurcación} (aquí, silla-nodo; capítulo~6).

En dimensión tres o más, un sistema determinista, sin nada aleatorio, puede tener soluciones acotadas, no periódicas y con dependencia sensible de las condiciones iniciales. El ejemplo canónico es el de Lorenz \citep{Lor63},
\begin{equation}
  \dot x=\sigma(y-x),\qquad \dot y=x(\rho-z)-y,\qquad \dot z=xy-\beta z,
  \label{eq:in-lorenz}
\end{equation}
con $\sigma=10$, $\rho=28$, $\beta=8/3$. Dos condiciones iniciales a $10^{-9}$ una de la otra son indistinguibles durante unos 25 segundos; hacia $t\approx33$ la diferencia ya es del tamaño del atractor (\cref{fig:in-panorama}f). \clave{Determinismo no implica predictibilidad.} Este no es un curso de caos, pero conviene saber dónde termina el mapa.

\video{V1_1_lorenz}{Dos trayectorias de Lorenz a $10^{-9}$}
\codigorepo{cap01_introduccion/C1_Ej3_Panorama.m}{lorenz}{Sensibilidad a las condiciones iniciales}

\begin{trampa}[label=tr:in-caos]{Una simulación muy precisa predice el futuro}
\plausible{Con tolerancias de $10^{-12}$, la simulación de Lorenz es la trayectoria verdadera hasta $t=40$.}
Un error de $10^{-12}$ en cada paso se amplifica como la diferencia de condiciones iniciales: en promedio, por un factor $e^{0.9}$ por unidad de tiempo. Hacia $t\approx30$ el error de redondeo ya es del tamaño del atractor, y el instante exacto en que la simulación se separa de la solución verdadera cambia si se cambian las tolerancias (pruébelo en \cmd{C1\_Ej3\_Panorama}).
\correcto{En un sistema caótico la simulación es confiable en un horizonte finito y, después, solo en lo cualitativo: la forma del atractor, las estadísticas, el rango de valores. Ningún integrador predice $x(40)$.}
\end{trampa}

\begin{figure}[ht]
\centering
\includegraphics[width=\linewidth]{cap01/t_mapa_curso}
\caption{El mapa del curso. Los capítulos 1--6 construyen la geometría del espacio de estado; los capítulos 7--10, la teoría de estabilidad que no necesita resolver la ecuación; los capítulos 11 y 12, el diseño.}
\label{fig:in-mapa}
\end{figure}

\begin{table}[ht]
\centering\small
\begin{tabular}{@{}p{0.42\linewidth}p{0.5\linewidth}@{}}
\toprule
Fenómeno & Dónde se estudia\\
\midrule
Equilibrios, estabilidad en la recta y en el plano & capítulos~\ref{cap:unidimensionales} y~4\\
Existencia, unicidad, escape en tiempo finito & capítulos~\ref{cap:unidimensionales} y~\ref{cap:preliminares}\\
Órbitas, ciclos límite, sincronización & \cref{cap:orbitas}\\
Bifurcaciones & \cref{cap:bifurcaciones}\\
Estabilidad de Lyapunov, región de atracción, LaSalle & capítulos~\ref{cap:lyapunov} y~8\\
Sistemas variantes en el tiempo, ISS, CLF & capítulos~9 y~10\\
Diseño: pasividad y linealización por retroalimentación & \cref{cap:pasividad,cap:linealizacion}\\
\bottomrule
\end{tabular}
\caption{Dónde se estudia cada fenómeno.}
\label{tab:in-mapa}
\end{table}

%% ---------------------------------------------------------------------
\section{Un ejemplo con historia: el crecimiento de una población}
\label{sec:in-poblaciones}
%% ---------------------------------------------------------------------

Volvamos a los conejos del \cref{ej:in-islas}. El modelo $P[k+1]=2P[k]$ duplica la población cada año sin límite: en unos cientos de generaciones los conejos cubrirían el planeta. Hace falta un freno.

\begin{historia}[Malthus, Verhulst y la ecuación logística]
Thomas Robert Malthus (1766--1834), pastor anglicano y después profesor de economía política en el colegio de la Compañía de las Indias Orientales en Haileybury, publicó en 1798 su \emph{Ensayo sobre el principio de la población} \citep{Mal86}: la población tiende a crecer en progresión geométrica y los alimentos solo en progresión aritmética, de modo que los recursos acaban por limitar a la población. Pierre François Verhulst (1804--1849), matemático belga doctorado en Gante y discípulo de Adolphe Quetelet, convirtió esa idea en una ecuación: en 1838 propuso un crecimiento que se frena al acercarse a un máximo, y en 1845 la llamó \emph{logística} \citep{Ver45}. La ecuación quedó olvidada hasta que Raymond Pearl y Lowell Reed la redescubrieron en 1920 para ajustar los censos de Estados Unidos \citep{PeRe20}. En 1976, Robert May (1936--2020) mostró en \emph{Nature} que su versión \emph{discreta}, la más simple que uno pueda imaginar, tiene una dinámica extraordinariamente complicada \citep{May76}; a partir de ese artículo los biólogos empezaron a tomarse en serio los modelos no lineales, y dos años después Mitchell Feigenbaum encontró en ella constantes universales del camino al caos \citep{Fei78}.
\end{historia}

\subsection{El mapa logístico}

Con $P[k]$ la población en el año $k$, $K$ la capacidad del medio y $R$ el factor de crecimiento por generación,
\begin{equation}
  P[k+1]=R\,P[k]\Bigl(1-\frac{P[k]}{K}\Bigr).
  \label{eq:in-mapa}
\end{equation}
El factor $(1-P/K)$ es la sobrepoblación: cuando $P$ se acerca a $K$, las crías mueren antes de reproducirse. Los equilibrios (puntos fijos, $P^*=F(P^*)$) son $P^*=0$ y $P^*=K(1-1/R)$.

\begin{ejemplo}[label=ej:in-mapa]{Iterar el mapa}
Con natalidad $2$, mortandad $0.4$ (de modo que $R=1.6$), capacidad $K=32$ y $P[0]=1$: ¿cuánto vale $P[1]$? ¿Y $P[21]$? Repita con $R=3.2$, $3.5$ y $3.9$. ¿Puede un modelo lineal de una sola variable hacer lo mismo?
\solucion $P[1]=1.6\cdot1\cdot\tfrac{31}{32}=1.55$. Iterando, la población sube monótonamente hacia el punto fijo $K(1-1/R)=32\cdot0.375=12$; en $P[21]$ ya vale $12.0000$ con cuatro cifras. Con $R=3.2$ la población termina alternando entre $16.42$ y $25.58$: un \emph{ciclo de período 2}. Con $R=3.5$ el ciclo es de período 4. Con $R=3.9$ no se repite nunca (\cref{fig:in-logistica}). Un modelo lineal escalar $P[k+1]=aP[k]$ solo puede crecer, decrecer o alternar de signo geométricamente: nada de eso.
\end{ejemplo}

\begin{figure}[ht]
\centering
\includegraphics[width=\linewidth]{cap01/f_logistica}
\caption{El mapa logístico \eqref{eq:in-mapa} con $K=32$, $P[0]=1$ y cuatro valores de $R$.}
\label{fig:in-logistica}
\end{figure}

\codigorepo{cap01_introduccion/C1_Ej2_Logistica.m}{discreta}{Iterar el mapa logístico}
\codigorepo{cap01_introduccion/C1_Ej2_Logistica.m}{regimenes}{Puntos fijos, ciclos y caos}

¿Por qué aparece el ciclo de período 2? El punto fijo $P^*=K(1-1/R)$ es estable si $\abs{F'(P^*)}<1$, y $F'(P^*)=2-R$. Para $1<R<3$ se cumple; en $R=3$, $F'(P^*)=-1$ y el punto fijo pierde la estabilidad: nace un ciclo de período 2. Al seguir subiendo $R$ ese ciclo también se desestabiliza y nace uno de período 4, luego de 8, \dots, y en $R\approx3.57$ empieza el caos. El \clave{diagrama de bifurcación} de la \cref{fig:in-bifurcacion} resume ese camino.

\video{V1_2_logistica}{El mapa logístico, de la calma al caos}

\subsection{La misma ecuación en tiempo continuo}

La versión continua es la ecuación de Verhulst,
\begin{equation}
  \dot P=R\,P\Bigl(1-\frac PK\Bigr),\qquad P(0)=P_0,
  \label{eq:in-logistica}
\end{equation}
que tiene solución explícita (verifíquelo derivando):
\[
  P(t)=\frac{K\,P_0\,e^{Rt}}{K+P_0\,(e^{Rt}-1)} .
\]
Sus equilibrios son $0$ y $K$, y para todo $R>0$ y $P_0\in(0,K)$ la solución sube monótonamente hasta $K$. El crecimiento $\dot P$ es máximo cuando $P=K/2$; la tasa \emph{per cápita}, $\dot P/P=R(1-P/K)$, es una recta que vale $R$ en $P=0$ y se anula en $K$. En el \cref{cap:unidimensionales} se verá por qué una ecuación escalar continua como \eqref{eq:in-logistica} no puede oscilar, ni mucho menos ser caótica.

\begin{figure}[ht]
\centering
\includegraphics[width=\linewidth]{cap01/f_bifurcacion}
\caption{Izquierda: diagrama de bifurcación del mapa logístico; las líneas marcan los cuatro valores de la \cref{fig:in-logistica}. Derecha: la ecuación logística continua con los mismos $R$ siempre converge monótonamente a $K$.}
\label{fig:in-bifurcacion}
\end{figure}

\codigorepo{cap01_introduccion/C1_Ej2_Logistica.m}{continua}{La logística continua y su solución exacta}

\begin{trampa}[label=tr:in-discreta]{La versión discreta y la continua describen lo mismo}
\plausible{\eqref{eq:in-mapa} y \eqref{eq:in-logistica} son ``la misma ecuación logística'': deben comportarse igual.}
Una es dócil (dos equilibrios, monotonía, solución explícita) y la otra es salvaje (ciclos de todos los períodos y caos). Ni siquiera coinciden los equilibrios: $K(1-1/R)$ contra $K$. En el mapa, $R$ es un \emph{factor por generación}; en la ecuación diferencial es una \emph{tasa} (por unidad de tiempo). Y en tiempo discreto la solución puede ``saltar'' por encima del equilibrio, algo imposible para una solución continua de una ecuación escalar.
\correcto{Discretizar no es una operación inocente. Un mapa y una ecuación diferencial ``parecidos'' pueden tener dinámicas cualitativamente distintas, y lo mismo vale para un método numérico con paso demasiado grande.}
\end{trampa}

%% ---------------------------------------------------------------------
\section{La misma matemática con distintas etiquetas}
\label{sec:in-energia}
%% ---------------------------------------------------------------------

Un argumento fuerte para estudiar los sistemas no lineales en abstracto es que un puñado de estructuras reaparece, con otros nombres, en campos que no se hablan entre sí. Durante el semestre usted trabajará sobre un sistema que puede venir de la biología, de los circuitos, de los procesos o de la mecánica; las herramientas son las mismas. Esta sección muestra dos de esas estructuras: la energía, que anticipa el método de Lyapunov, y un sistema que se puede controlar pero no estabilizar.

%% ---------------------------------------------------------------------
\subsection{La energía: concluir sin resolver}
%% ---------------------------------------------------------------------

Un bloque de masa $m$ unido a una pared por un resorte de constante $k$, con fricción viscosa $c$ y una fuerza $u$, cumple $m\ddot y=u-c\dot y-ky$. Con $x_1=y$ y $x_2=\dot y$,
\begin{equation}
  \dot x_1=x_2,\qquad \dot x_2=-\frac km\,x_1-\frac cm\,x_2+\frac1m\,u .
  \label{eq:in-masa}
\end{equation}
Con $u=0$, considere la energía total, potencial más cinética,
\begin{equation}
  E(x)=\tfrac12k\,x_1^2+\tfrac12m\,x_2^2 .
  \label{eq:in-energia}
\end{equation}
¿Cómo cambia a lo largo de una solución? Por la regla de la cadena,
\begin{equation}
  \frac{\dif}{\dif t}E\bigl(x(t)\bigr)=\inner{\nabla E(x)}{f(x)}=kx_1x_2+mx_2\Bigl(-\frac kmx_1-\frac cmx_2\Bigr)=-c\,x_2^2\le0 .
  \label{eq:in-dE}
\end{equation}
Sabemos dos cosas: $E\ge0$, por ser suma de cuadrados, y $E$ no crece mientras el bloque se mueve. Luego la energía tiene límite y el bloque tiene que ir dejando de moverse. Y lo dedujimos \alerta{sin resolver la ecuación diferencial}. Esa es la idea que Lyapunov generalizó en 1892, reemplazando ``energía'' por ``cualquier función que se comporte como una energía'', y es el núcleo del curso (capítulo~7).

\begin{figure}[ht]
\centering
\includegraphics[width=\linewidth]{cap01/f_energia}
\caption{Masa-resorte con $m=1$, $k=4$, $c=0.5$ desde $x(0)=(1,0)$. La energía nunca crece, pero se aplana en los instantes en que $x_2=0$ (punteado), donde $\dot E=0$.}
\label{fig:in-energia}
\end{figure}

\codigorepo{cap01_introduccion/C1_Ej4_Energia.m}{energia}{La energía a lo largo de una solución}

\begin{trampa}[label=tr:in-lasalle]{La energía decrece, luego el sistema va al origen}
\plausible{Como $\dot E=-cx_2^2\le0$, el bloque se detiene en $x=0$.}
La desigualdad dice que $E$ no crece; no dice que llegue a cero, y $\dot E$ se anula en toda la recta $x_2=0$, no solo en el origen (\cref{fig:in-energia}): en cada inversión del movimiento el bloque está quieto y $\dot E=0$. Un $\dot E$ \emph{semidefinido} negativo no basta para concluir que $x_1\to0$; de ahí solo se sigue, con un argumento extra, que $x_2\to0$.
\correcto{Hace falta un paso más: una solución que se quedara en $\{x_2=0\}$ tendría $\dot x_2=0$, luego $x_1=0$. Ninguna solución puede vivir en $\{\dot E=0\}$ salvo el origen, y eso, hecho preciso, es el \clave{principio de invarianza de LaSalle} (capítulo~8). Saltarse ese paso es la trampa más común del curso.}
\end{trampa}

Para el péndulo \eqref{eq:in-pendulo}, la energía $E=mg\ell(1-\cos x_1)+\tfrac12m\ell^2x_2^2$ cumple $\dot E=-k\ell^2x_2^2\le0$. Ahora hay una complicación: $E$ es positiva cerca del origen, pero se anula también en $x_1=2\pi$, $4\pi$, \dots\ (el péndulo da una vuelta y vuelve a colgar). La energía sirve como argumento \emph{local}, en una vecindad del equilibrio de abajo, y esa restricción es el germen de la estabilidad local y de la región de atracción.

\paragraph{Otras etiquetas, la misma estructura.} El péndulo sobre un carro modela una grúa (péndulo colgando), un cohete o un \emph{segway} (péndulo invertido): son el mismo sistema mecánico $M(q)\ddot q+C(q,\dot q)+G(q)=B(q)u$ con equilibrios distintos \citep{KeBr23}. Para escribirlo en la forma \eqref{eq:in-general} hay que invertir la matriz de inercia $M(q)$, y verificar que es invertible en todo el espacio de estado: la forma $\dot x=f(x,u)$ no siempre es gratis.

%% ---------------------------------------------------------------------
\subsection{Controlable, pero no estabilizable}
\label{sec:in-brockett}
%% ---------------------------------------------------------------------

Un robot móvil de dos ruedas (el \emph{uniciclo}), con posición $(p_1,p_2)$, orientación $\phi$, velocidad de avance $u_1$ y velocidad de giro $u_2$, es
\[
  \dot p_1=u_1\cos\phi,\qquad \dot p_2=u_1\sin\phi,\qquad\dot\phi=u_2 .
\]
Con un cambio de coordenadas y de entradas se lleva a la forma del \clave{integrador no holonómico} o \clave{integrador de Brockett} \citep{Bro83,KeBr23}:
\begin{equation}
  \dot z_1=v_1,\qquad \dot z_2=v_2,\qquad \dot z_3=z_1v_2-z_2v_1 .
  \label{eq:in-brockett}
\end{equation}

\emph{Es controlable.} La velocidad instantánea nunca apunta en la dirección $z_3$ desde el origen, pero dando una vuelta sí se llega. Con $v=\rho(-\sin t,\cos t)$ desde $z=0$, $(z_1,z_2)=\rho(\cos t-1,\sin t)$ recorre un círculo y regresa al origen en $t=2\pi$, mientras que
\[
  \dot z_3=z_1v_2-z_2v_1=\rho^2(1-\cos t)\quad\Longrightarrow\quad z_3(2\pi)=2\pi\rho^2 ,
\]
el doble del área encerrada (\cref{fig:in-brockett}). Combinando vueltas se llega de cualquier punto a cualquier otro. Note que la linealización en el origen, $\dot z_1=v_1$, $\dot z_2=v_2$, $\dot z_3=0$, \emph{no} es controlable: su matriz de controlabilidad tiene rango $2<3$.

\begin{figure}[ht]
\centering
\includegraphics[width=\linewidth]{cap01/f_brockett}
\caption{Integrador de Brockett: una vuelta de radio $\rho=0.5$ en $(z_1,z_2)$ deja $(z_1,z_2)$ donde estaba y sube $z_3$ en $2\pi\rho^2=\pi/2$.}
\label{fig:in-brockett}
\end{figure}

\codigorepo{cap01_introduccion/C1_Ej8_Brockett.m}{vuelta}{Una vuelta sube $z_3$}

\emph{Y no es estabilizable con una realimentación continua.} Brockett encontró una condición necesaria:

\begin{teorema}[label=thm:in-brockett]{Condición de Brockett \citep{Bro83}}
Sea $\dot x=f(x,u)$ con $f$ de clase $\Cont1$ y $f(0,0)=0$. Si existe una realimentación $u=k(x)$, de clase $\Cont1$ y con $k(0)=0$, que hace asintóticamente estable el origen del lazo cerrado, entonces la imagen de $f$ cerca del origen contiene una vecindad del origen: para todo $\varepsilon>0$ existe $\delta>0$ tal que todo $y\in\R^n$ con $\norm y<\delta$ se puede escribir como $y=f(x,u)$ con $\norm x<\varepsilon$ y $\norm u<\varepsilon$.
\end{teorema}

\notamargen{Brockett lo enunció con $k$ suave; Zabczyk (1989) lo extendió a realimentaciones continuas con soluciones únicas. La prueba usa topología (grado de un mapa) y está fuera del alcance de estas notas.}
Para \eqref{eq:in-brockett}, tome $y=(0,0,\delta)$: las dos primeras componentes de $y=f(z,v)$ obligan a $v_1=v_2=0$, y entonces la tercera es $z_1\cdot0-z_2\cdot0=0\neq\delta$. Ningún punto de la forma $(0,0,\delta)$, por pequeño que sea $\delta$, es una velocidad posible. La condición falla, y \alerta{ninguna realimentación estática y suave estabiliza el origen}, aunque el sistema sea controlable.

\codigorepo{cap01_introduccion/C1_Ej8_Brockett.m}{brockett}{Falla la condición de Brockett}

\begin{trampa}[label=tr:in-controlable]{Controlable implica estabilizable}
\plausible{El integrador de Brockett es controlable; luego hay una realimentación de estado que lo estabiliza.}
En sistemas lineales eso es cierto: si $(A,B)$ es controlable, se ubican los polos donde se quiera con $u=-Kx$. En sistemas no lineales el puente se rompe: el \cref{thm:in-brockett} muestra que el integrador de Brockett, controlable, no se estabiliza con ninguna $u=k(x)$ suave.
\correcto{No es que no hayamos encontrado el controlador: no existe. Para estabilizar un robot móvil hay que usar leyes discontinuas, variantes en el tiempo o híbridas. Este es, quizás, el mejor argumento de por qué hace falta una teoría no lineal.}
\end{trampa}

%% ---------------------------------------------------------------------
\section{¿Qué aportan los sistemas no lineales al control?}
\label{sec:in-control}
%% ---------------------------------------------------------------------

Un controlador se usa para seguir una referencia, regular un estado, garantizar estabilidad, dar un transitorio adecuado y reducir la sensibilidad a los parámetros de la planta. Muchas veces un controlador lineal diseñado sobre la linealización basta. Esta sección muestra tres situaciones en las que no: una en la que el lazo lineal no funciona, otra en la que una ley no lineal funciona mejor, y la del \cref{thm:in-brockett}, en la que ninguna ley continua funciona.

%% ---------------------------------------------------------------------
\subsection{El lazo y lo que se le olvida al diagrama de bloques}
%% ---------------------------------------------------------------------

En \clave{lazo abierto} la entrada se calcula sin mirar la planta: una tostadora calienta el pan un tiempo fijo, cualquiera que sea la temperatura que alcance. En \clave{lazo cerrado} la entrada depende del error entre lo que se mide y lo que se quiere, $e=y_m-r$: un aire acondicionado. El controlador se \emph{diseña} con un modelo, pero se \emph{evalúa} sobre la planta real, donde hay perturbaciones $v$ en la entrada y errores $w$ en la medición (\cref{fig:in-lazo}). Si además no se mide todo el estado, hace falta un \clave{observador} que lo estime.

\begin{figure}[ht]
\centering
\includegraphics[width=0.95\linewidth]{cap01/t_lazo}
\caption{El lazo de control con perturbación de entrada $v$ y error de medición $w$. El controlador solo ve $y_m$.}
\label{fig:in-lazo}
\end{figure}

\begin{ejemplo}[label=ej:in-pid]{Lo que el integrador corrige y lo que no}
Para el masa-resorte $m\ddot y+c\dot y+ky=u+v$ con $v$ y $w$ constantes y desconocidos, compare el régimen permanente con un PD, $u=-k_Pe-k_D\dot e$, y con un PID, que agrega $-k_I\int e$.
\solucion Con el PD, en el equilibrio $\dot y=0$ y $0=-k_P(y+w-r)+v-ky$, de donde
\[
  y^*=\frac{k_P(r-w)+v}{k+k_P} .
\]
Con $m=1$, $c=1$, $k=2$, $k_P=10$, $k_D=3$, $r=1$, $v=0.5$, $w=0.1$: $y^*=0.7917$. Con el PID, el estado del integrador solo deja de cambiar si $e=0$, es decir, $y^*+w=r$: $y^*=r-w=0.9$ (\cref{fig:in-pid}). El integrador elimina el efecto de la perturbación de entrada (el \clave{principio del modelo interno}) pero no el del sensor: el controlador hace que la \emph{medición} valga $r$, y la medición miente.
\end{ejemplo}

\begin{figure}[ht]
\centering
\includegraphics[width=0.78\linewidth]{cap01/f_pid}
\caption{PD y PID con perturbación $v=0.5$ y error de sensor $w=0.1$. Ninguno llega a la referencia.}
\label{fig:in-pid}
\end{figure}

\codigorepo{cap01_introduccion/C1_Ej9_PID.m}{pid}{PD, PID, perturbación y ruido}

%% ---------------------------------------------------------------------
\subsection{Cuando la ley lineal no funciona}
%% ---------------------------------------------------------------------

\begin{ejemplo}[label=ej:in-cubico]{La planta $\dot x=x+u^3$ \citep{Mar03}}
Considere $\dot x=x+u^3$, $x(0)=1$. Compare $u=0$, la realimentación lineal $u=-kx$ y la no lineal $u=-(kx)^{1/3}$, con $k=5$.
\solucion Con $u=0$, $x(t)=e^t$: el origen es inestable. Con $u=-kx$ el lazo cerrado es
\[
  \dot x=x-k^3x^3 ,
\]
que tiene \emph{tres} equilibrios, $0$ y $\pm k^{-3/2}=\pm0.0894$. La derivada en el origen es $1$: \alerta{el origen sigue siendo inestable}, para cualquier $k$, y la solución desde $x(0)=1$ converge a $0.0894$, no a $0$. La realimentación lineal no estabilizó nada: creó dos equilibrios nuevos (\cref{fig:in-cubico}). Con $u=-(kx)^{1/3}$ (la raíz cúbica real), $u^3=-kx$ y el lazo cerrado es $\dot x=(1-k)x=-4x$: exponencialmente estable.
\end{ejemplo}

\begin{figure}[ht]
\centering
\includegraphics[width=\linewidth]{cap01/f_control_cubico}
\caption{Izquierda: $x(t)$ con las tres leyes. Derecha: el campo en lazo cerrado. Con $u=-kx$ el origen (círculo vacío) sigue siendo inestable y aparecen dos equilibrios estables en $\pm k^{-3/2}$.}
\label{fig:in-cubico}
\end{figure}

\codigorepo{cap01_introduccion/C1_Ej6_ControlCubico.m}{leyes}{Tres leyes para $\dot x=x+u^3$}

\begin{trampa}[label=tr:in-cubico]{Converge a algo, luego es estable}
\plausible{Con $u=-kx$ la solución deja de crecer y se queda en un valor constante pequeño: el origen es estable, aunque no asintóticamente.}
Así lo decían las versiones anteriores de estas notas. Pero la solución no converge al origen sino a \emph{otro} equilibrio, $k^{-3/2}$, y el origen es inestable: toda condición inicial distinta de cero, por pequeña que sea, se aleja de él. ``Estable'' es una propiedad de un equilibrio, y aquí hay tres.
\correcto{Antes de interpretar una simulación, calcule todos los equilibrios del lazo cerrado y su estabilidad. La razón de fondo: la linealización de la planta en el origen, $\dot x=x+0\cdot u$, no depende de $u$; ninguna ley lineal puede moverle el polo.}
\end{trampa}

%% ---------------------------------------------------------------------
\subsection{Cuando la ley no lineal funciona mejor}
%% ---------------------------------------------------------------------

\begin{ejemplo}[label=ej:in-vsc]{Doble integrador y estructura variable}
Considere el doble integrador $\dot x_1=x_2$, $\dot x_2=u$, desde $x(0)=(1,1)$. Compare $u=-kx_1$, el PD $u=-x_1-x_2$ y la ley de \clave{estructura variable}
\begin{equation}
  u=k_3\,x_1,\qquad k_3=\begin{cases}+1 & \text{si } x_1s<0,\\ -1 & \text{si } x_1s>0,\end{cases}\qquad s=x_1+x_2 .
  \label{eq:in-vsc}
\end{equation}
\solucion Con $u=-kx_1$ el lazo es un oscilador armónico: acotado, pero sin converger. El PD da polos en $-\tfrac12\pm\tfrac{\sqrt3}2\,j$: converge, con un sobrepico de $-0.21$, y tarda $8.91$ en quedar dentro del 2\,\% de la norma inicial. La ley \eqref{eq:in-vsc} alterna dos sistemas que, cada uno por separado, \emph{no} convergen: con $k_3=-1$ un centro ($\ddot x_1=-x_1$) y con $k_3=+1$ una silla ($\ddot x_1=x_1$). Desde $(1,1)$, $x_1s>0$ y el estado gira por el círculo $x_1=\cos t+\sin t$, $x_2=\cos t-\sin t$ hasta tocar la recta $s=0$ en $t=\pi/2$, en el punto $(1,-1)$. Allí cada sistema empuja el estado de vuelta hacia la recta, y el estado \emph{desliza} por ella: sobre $s=0$, $\dot x_1=-x_1$. Llega al 2\,\% en $\pi/2+\ln50=5.48$ y sin sobrepico (\cref{fig:in-vsc}).
\end{ejemplo}

\begin{figure}[ht]
\centering
\includegraphics[width=\linewidth]{cap01/f_doble_integrador}
\caption{Doble integrador con tres leyes. La de estructura variable llega a la recta $s=0$ en $t=\pi/2$ y desliza por ella.}
\label{fig:in-vsc}
\end{figure}

\codigorepo{cap01_introduccion/C1_Ej7_DobleIntegrador.m}{vsc}{Cuatro leyes para el doble integrador}

\begin{modelo}[EstructuraVariable.slx]
\cmd{crearModelo\_EstructuraVariable.m} arma el lazo con dos integradores y un bloque \emph{Switch} que elige entre $+x_1$ y $-x_1$ según el signo de $x_1s$, simula 10\,s y compara con \cmd{ode45}.
\end{modelo}

\begin{trampa}[label=tr:in-discontinuo]{El lazo con conmutación es una ecuación diferencial como cualquier otra}
\plausible{La ley \eqref{eq:in-vsc} da un lazo $\dot x=f(x)$; el teorema de existencia y unicidad garantiza que la simulación es la solución.}
Sobre la recta $s=0$ el lado derecho es \emph{discontinuo}, y los teoremas de existencia y unicidad del \cref{cap:preliminares} no aplican. ¿Qué significa ``solución'' cuando el estado desliza por una superficie donde $f$ no está definida? Filippov dio la respuesta: una solución en el sentido de las inclusiones diferenciales \citep{Fil88}. En la simulación, \cmd{ode45} conmuta muchas veces por paso (\emph{chattering}) y se vuelve lento.
\correcto{Los lazos con conmutación son útiles (el control por modos deslizantes, \citealp{Utk77}), pero exigen otra teoría. Y la comparación con el PD depende de las ganancias: con otras ganancias el PD es más rápido. Lo que no puede hacer ninguna ley lineal es llevar el estado a una recta y mantenerlo allí con dinámica de primer orden.}
\end{trampa}

Los dos primeros ejemplos muestran ventajas de desempeño; el de la \cref{sec:in-brockett} muestra algo cualitativamente distinto: una \emph{imposibilidad}. El control no lineal no es ``control lineal, pero mejor''. A veces es la única opción, y a veces es la teoría que dice que lo que se busca no existe.

%% ---------------------------------------------------------------------
\section{Un poco de historia}
\label{sec:in-historia}
%% ---------------------------------------------------------------------

Las herramientas del curso tienen dos padres, Henri Poincaré y Aleksandr Lyapunov, y una historia que vale la pena conocer porque explica \emph{qué preguntas} se hacen sobre un sistema no lineal.

\paragraph{De Newton a Poincaré: cambiar la pregunta.} Newton inventó las ecuaciones diferenciales y resolvió el problema de dos cuerpos: la órbita de un planeta alrededor del Sol. Durante dos siglos los matemáticos intentaron hacer lo mismo con tres cuerpos (Sol, Tierra y Luna) buscando fórmulas, sin éxito. Lagrange enunció en su \emph{Mecánica analítica} (1788) que un mínimo aislado de la energía potencial de un sistema conservativo es un equilibrio estable; la demostración rigurosa la dio Dirichlet en 1846. Poincaré cambió la pregunta: en lugar de calcular dónde estará un planeta en cada instante, preguntó si el sistema solar es estable, si los planetas se quedarán para siempre en una región acotada o alguno escapará.
\notamargen{Este cambio de pregunta, de la cantidad a la cualidad, es la lección metodológica más importante del capítulo: buena parte del semestre consiste en aprender qué preguntas \emph{sí} se pueden responder sobre un sistema no lineal.}

\begin{historia}[Henri Poincaré (1854--1912)]
Nació en Nancy el 29 de abril de 1854 y murió en París el 17 de julio de 1912. Estudió en la Escuela Politécnica (1873--1875) y en la Escuela de Minas, de la que salió ingeniero en 1879; ese mismo año presentó su tesis doctoral sobre ecuaciones en derivadas parciales. Entre 1881 y 1886 publicó la \emph{Memoria sobre las curvas definidas por una ecuación diferencial} \citep{Poi1881}, donde nace la \clave{teoría cualitativa} de las ecuaciones diferenciales: la clasificación de los puntos singulares (nodos, sillas, focos, centros), el índice, los ciclos límite y la idea de seguir una trayectoria hasta que vuelve a cortar una curva (el mapa de Poincaré). En 1889 ganó el premio del rey Óscar~II de Suecia con una memoria sobre el problema de los tres cuerpos. Una pregunta de Edvard Phragmén, que preparaba la memoria para su publicación, llevó a Poincaré a encontrar en ella un error grave; al corregirlo descubrió que las trayectorias cerca de ciertas órbitas se enredan de manera inextricable: la primera visión de lo que hoy llamamos caos \citep{Poi1890}. Pagó de su bolsillo, más de lo que valía el premio, la reimpresión de la versión corregida.
\end{historia}

\begin{historia}[Aleksandr Mijáilovich Lyapunov (1857--1918)]
Nació en Yaroslavl el 6 de junio de 1857. Estudió en San Petersburgo con Pafnuti Chebyshev, quien le planteó una pregunta sobre la estabilidad de las formas de equilibrio de un fluido que rota; de ahí salió su tesis de maestría (1884). En 1892 presentó en la Universidad de Moscú su tesis doctoral, \emph{El problema general de la estabilidad del movimiento}, escrita mientras era profesor en Járkov. Allí aparecen los dos métodos que usamos: el indirecto (la linealización) y el directo, que concluye la estabilidad con una función que se comporta como una energía, sin resolver la ecuación. La obra se tradujo al francés en 1907 y al inglés solo en 1992 \citep{Lya92}, lo que explica por qué el método directo tardó medio siglo en llegar a la ingeniería occidental. Murió en Odesa el 3 de noviembre de 1918, días después de la muerte de su esposa \citep{MBM07}.
\end{historia}

\paragraph{Osciladores y servomecanismos (1920--1960).} Balthasar van der Pol estudió en 1926 las oscilaciones de relajación de los circuitos con tríodos \citep{van26}, y Alfred Liénard demostró en 1928 la existencia de una solución periódica única para una familia de osciladores \citep{Lie28}. En 1929, Aleksandr Andronov mostró que las oscilaciones autosostenidas de la radiotecnia son los ciclos límite de Poincaré \citep{And29}, y la escuela soviética (Mandelstam, Andronov, Vitt, Jaikin) construyó sobre esa idea la teoría de las oscilaciones no lineales. Georg Duffing publicó en 1918 su estudio de las oscilaciones forzadas con un resorte cúbico \citep{Duf18}. La Segunda Guerra Mundial impulsó el control de servomecanismos (control de tiro, guiado) y con él tres métodos: el plano de fase, la función descriptiva (que solo balancea el primer armónico) y el análisis de sistemas con relés. En 1944, Lur'e y Postnikov plantearon el problema de la estabilidad absoluta \citep{LuPo44}.

\paragraph{La era moderna (desde 1960).} En 1960, el año del primer congreso de la IFAC, Kalman y Bertram llevaron el método directo de Lyapunov al control \citep{KaBe60}, y LaSalle publicó su principio de invarianza \citep{LaS60}, precedido por Barbashin y Krasovskii \citep{BaKr52}. Popov dio en 1961 su criterio de estabilidad absoluta \citep{Pop61}; Yakubovich (1962) y Kalman (1963) lo conectaron con la teoría de estado en el lema KYP \citep{Yak62,Kal63}; Sandberg (1964) y Zames (1966) formularon el teorema de la ganancia pequeña y el criterio del círculo \citep{San64,Zam66a}. Edward Lorenz encontró en 1963 un atractor extraño en un modelo de la atmósfera \citep{Lor63}. Jan Willems propuso en 1972 la teoría de los sistemas disipativos \citep{Wil72}, de donde sale la pasividad (\cref{cap:pasividad}).

En los años ochenta y noventa el control no lineal pasó de describir a \emph{construir} \citep{KoAr01}. Brockett mostró en 1983 los límites de la estabilización continua (\cref{thm:in-brockett}); Artstein introdujo las funciones de Lyapunov de control \citep{Ars83} y Sontag dio una fórmula universal para construir el controlador a partir de ellas \citep{Son89}; Sontag también introdujo la estabilidad entrada-estado (ISS) \citep{Son89a}, cuyas caracterizaciones completó con Wang \citep{SoWa95}. Byrnes e Isidori desarrollaron la dinámica cero y la estabilización de sistemas de fase mínima \citep{ByIs91,Isi95}, base de la linealización por retroalimentación (\cref{cap:linealizacion}), y Krstić, Kanellakopoulos y Kokotović consolidaron el \emph{backstepping} \citep{KKK95}.

\begin{table}[ht]
\centering\small
\begin{tabular}{@{}lp{0.78\linewidth}@{}}
\toprule
Año & Hito\\
\midrule
1788 & Lagrange: mínimo de la energía potencial y estabilidad (prueba de Dirichlet, 1846)\\
1838 & Verhulst: la ecuación logística\\
1881--86 & Poincaré: teoría cualitativa de las ecuaciones diferenciales\\
1892 & Lyapunov: el problema general de la estabilidad del movimiento\\
1926--29 & van der Pol, Liénard y Andronov: oscilaciones autosostenidas y ciclos límite\\
1944 & Lur'e y Postnikov: estabilidad absoluta\\
1960 & Kalman y Bertram: Lyapunov en control; LaSalle: principio de invarianza\\
1961--66 & Popov; Yakubovich y Kalman (KYP); Sandberg y Zames (ganancia pequeña, círculo)\\
1963 & Lorenz: caos determinista\\
1972 & Willems: sistemas disipativos\\
1975--78 & Kuramoto: sincronización; May: el mapa logístico; Feigenbaum: universalidad\\
1983 & Brockett: condición necesaria; Artstein: funciones de Lyapunov de control\\
1989 & Sontag: fórmula universal e ISS\\
1985--95 & Byrnes e Isidori: dinámica cero y control geométrico; \emph{backstepping}\\
\bottomrule
\end{tabular}
\caption{Algunos hitos. Más detalles en \citet{Ath96}, \citet{Mic96} y \citet{KoAr01}.}
\label{tab:in-historia}
\end{table}

%% ---------------------------------------------------------------------
\section{Cómo trabajar en este curso}
\label{sec:in-trabajo}
%% ---------------------------------------------------------------------

Esta sección no es sobre sistemas no lineales: es sobre cómo se trabaja con ellos durante el semestre, y por qué las notas están escritas como están.

\subsection{Tres niveles de análisis}

Todo análisis se puede leer buscando tres cosas, en orden creciente de dificultad.
\begin{enumerate}
  \item \clave{Correcto y completo.} Los equilibrios están todos, la linealización está bien calculada, las conclusiones se siguen de las premisas, la simulación corresponde a las ecuaciones.
  \item \clave{Consciente de sus límites.} Además, el análisis dice \emph{dónde deja de valer}: si la conclusión es local o global, en qué región se cumplen las hipótesis, qué pasa si un parámetro se sale del rango supuesto, si un resultado numérico es una propiedad del sistema o un artefacto del integrador.
  \item \clave{Original.} Además, contiene algo que no está en ningún libro ni se obtiene pidiéndoselo a una herramienta: una conexión inesperada entre dos sistemas, una interpretación física que ilumina el resultado, un contraejemplo propio.
\end{enumerate}
El segundo nivel es donde ocurre la mayor parte del aprendizaje y donde se concentran los errores más costosos.

\subsection{Conclusiones plausibles y falsas}

Un sistema no lineal casi siempre admite una conclusión que \emph{suena} correcta, la que sale de razonar por analogía con el caso lineal, y que es falsa. Este capítulo ya mostró varias: la estabilidad como propiedad del sistema (\cref{sec:in-fenomenos}); ``la energía decrece, luego va al origen'' (\cref{tr:in-lasalle}); ``converge a algo, luego es estable'' (\cref{tr:in-cubico}); ``controlable, luego estabilizable'' (\cref{tr:in-controlable}); ``la logística discreta y la continua son lo mismo'' (\cref{tr:in-discreta}); ``una simulación precisa predice el futuro'' (\cref{tr:in-caos}). Hay más en cada capítulo, en las cajas rojas marcadas \textbf{Trampa}.

\begin{ideaclave}
En todos los casos la afirmación falsa es una \emph{generalización indebida} de algo que sí es cierto en un contexto más restringido. Las herramientas automáticas de generación de texto son particularmente buenas produciendo exactamente ese tipo de afirmación, porque producen lo que suena bien en el contexto. Distinguir lo plausible de lo correcto es la habilidad que el curso quiere desarrollar, y no se puede delegar.
\end{ideaclave}

\subsection{Tres reglas para simular}

El curso usa MATLAB y Simulink de manera intensiva; cada figura de estas notas sale de un script del repositorio (\cmd{codigo/}), y \cmd{verificarCapNN} comprueba sus números.
\begin{enumerate}
  \item \clave{La simulación no demuestra; ilustra.} Cien condiciones iniciales que convergen no prueban estabilidad global; una que no converge la refuta.
  \item \clave{Antes de simular, sepa qué espera ver.} Calcule los equilibrios, linealice, estime las escalas de tiempo. Si la simulación contradice el cálculo, uno de los dos está mal, y averiguar cuál es donde está el aprendizaje.
  \item \clave{Desconfíe del integrador.} Sistemas rígidos, oscilaciones de relajación, soluciones que escapan y lados derechos discontinuos son los sospechosos habituales de resultados espurios. \cmd{ode45} no siempre es la respuesta.
\end{enumerate}

%% ---------------------------------------------------------------------
\section{Para llevar}
\label{sec:in-cierre}
%% ---------------------------------------------------------------------

\begin{parallevar}
\begin{enumerate}
  \item \textbf{Modelo:} $\dot x=f(t,x,u)$, $y=h(t,x,u)$; invariante o variante en el tiempo según $f$ dependa de $t$. Con $u=k(t,x)$ el lazo cerrado no tiene entrada: análisis y síntesis se encuentran.
  \item \textbf{Equilibrios:} $f(\eq)=0$. Trasladarlos al origen es gratis; reescalar reduce parámetros. Ambas cosas hay que deshacerlas al interpretar.
  \item \textbf{Lo que lo lineal no puede:} varios equilibrios aislados, ciclos límite, escape en tiempo finito, período que depende de la amplitud, caos. La estabilidad es de un equilibrio, no del sistema.
  \item \textbf{El todo no es la suma de las partes}, y la versión discreta de una ecuación no se comporta como la continua.
  \item \textbf{Energía:} si $E\ge0$ y $\dot E\le0$ se concluye algo sin resolver la ecuación; con $\dot E$ solo semidefinida hace falta LaSalle.
  \item \textbf{Control:} una ley lineal puede no estabilizar (y crear equilibrios nuevos); una no lineal puede hacerlo mejor; y hay sistemas controlables que ninguna ley continua estabiliza (Brockett).
\end{enumerate}
\end{parallevar}

\begin{trampa}[title={\faBug\hspace{0.45em}Las trampas del capítulo, en una caja}]{}
\begin{itemize}
  \item Agregar $t$ como estado no convierte un sistema variante en invariante (\cref{tr:in-tiempo}).
  \item Una simulación de un sistema caótico no predice el futuro lejano (\cref{tr:in-caos}).
  \item El mapa logístico y la ecuación logística no son lo mismo (\cref{tr:in-discreta}).
  \item $\dot E\le0$ no lleva al origen sin LaSalle (\cref{tr:in-lasalle}).
  \item Controlable no implica estabilizable (\cref{tr:in-controlable}).
  \item Converger a algo no es ser estable: puede ser otro equilibrio (\cref{tr:in-cubico}).
  \item Un lazo con conmutación no es una ecuación diferencial ordinaria (\cref{tr:in-discontinuo}).
\end{itemize}
\end{trampa}

\begin{profundiza}[Correcciones respecto a las versiones anteriores (Handout1, Handout1\_2026 e Intro2026)]
(i) En el ejemplo $\dot x=x+u^3$ con $u=-kx$ se decía que el origen quedaba ``estable pero no asintóticamente estable''. Es falso: el origen sigue siendo inestable y la solución converge a otro equilibrio, $k^{-3/2}$ (\cref{ej:in-cubico}). (ii) La tesis de Poincaré (1879) trató de ecuaciones en derivadas parciales; la clasificación de puntos singulares, el índice y los ciclos límite son de la \emph{Memoria} de 1881--1886. Nació el 29 de abril, no el 28. (iii) Duffing (1918) estudió oscilaciones forzadas con resorte cúbico; el caos en su ecuación se encontró mucho después. (iv) El criterio del círculo es de Sandberg y Zames; Popov dio en 1961 su propio criterio. (v) La dinámica cero es de Byrnes e Isidori; la ISS la introdujo Sontag en 1989, y con Wang (1995) dio sus caracterizaciones. (vi) La obra de Lyapunov se tradujo al francés 15 años después de 1892, no 25. (vii) La ecuación logística la redescubrieron Pearl y Reed en 1920; se eliminó la mención a un redescubrimiento por ``Alfred Volta''. Rousseau era amigo del padre de Malthus, no su mentor. Feigenbaum trabajó con el mapa logístico a finales de los setenta. (viii) Se corrigió la grafía de Artstein. (ix) La ley de estructura variable produce un lazo discontinuo: se agregó la advertencia sobre la noción de solución y se matizó la comparación con el PD. (x) Se quitaron los restos de la plantilla Tufte y la notación se unificó con \texttt{notacion.sty} (invariante/variante en el tiempo, $\eq$, $f(t,x,u)$). Todos los números se verificaron con \texttt{herramientas/cap01\_numeros.py} (36 de 36) y \cmd{verificarCap01.m} (34 de 34).
\end{profundiza}

%% ---------------------------------------------------------------------
\section{Problemas}
\label{sec:in-problemas}
%% ---------------------------------------------------------------------

Los problemas marcados con \etiquetaM{} requieren \matlab; los marcados con \etiquetaEstrella{} son de profundización. Cada uno tiene una trampa: búsquela antes de dar la respuesta.

\begin{problema}[equilibrios y fenómenos]\label{prob:in-equilibrios}
Encuentre todos los equilibrios y diga cuáles de los fenómenos de la \cref{sec:in-fenomenos} puede tener cada sistema: (a) $\dot x=x-x^3$; (b) $\dot x=x^2$; (c) $\dot x=\sin x$; (d) $\dot x=rx-x^2$, según el signo de $r$.
\end{problema}

\begin{problema}[Duffing]\label{prob:in-duffing}
Escriba en variables de estado la ecuación de Duffing $\ddot x+\delta\dot x+\alpha x+\beta x^3=\gamma\cos\omega t$. ¿Es invariante en el tiempo? Con $\gamma=0$, ¿cuántos equilibrios tiene según los signos de $\alpha$ y $\beta$?
\end{problema}

\begin{problema}[un circuito con diodo túnel]\label{prob:in-diodo}
Un circuito RLC tiene en paralelo con el condensador un diodo túnel de característica $i=h(v)$ no monótona (\cref{sec:un-diodo}). Escriba el modelo de estado con la corriente del inductor y el voltaje del condensador. ¿Cuántos equilibrios puede tener? Compare con el mismo circuito con una resistencia lineal.
\end{problema}

\begin{problema}[reescalar la logística]\label{prob:in-reescalar}
(a) Muestre que con $z=P/K$ y $\tau=Rt$ la ecuación \eqref{eq:in-logistica} queda $\dif z/\dif\tau=z(1-z)$, sin parámetros. (b) Haga lo mismo con la logística con cosecha $\dot P=RP(1-P/K)-H$. ¿Cuántos parámetros quedan, y cuál es? (Este sistema reaparece en el \cref{cap:unidimensionales}.)
\end{problema}

\begin{problema}[descubrir LaSalle]\label{prob:in-lasalle}
Para el masa-resorte \eqref{eq:in-masa} con $u=0$ y $c>0$, $\dot E=-cx_2^2$. Describa el conjunto donde $\dot E=0$ y argumente, sin teoremas, por qué la única solución que puede quedarse en ese conjunto para siempre es la del origen. ¿Qué cambia si $c=0$?
\end{problema}

\begin{problema}[la condición de Brockett]\label{prob:in-brockett}
(a) Verifique la condición del \cref{thm:in-brockett} para el doble integrador $\dot x_1=x_2$, $\dot x_2=u$: ¿todo $y$ pequeño es $f(x,u)$ para $x,u$ pequeños? (b) Repita con el integrador de Brockett. (c) ¿Qué concluye en cada caso? ¿Qué \emph{no} puede concluir del caso (a)?
\end{problema}

\begin{problema}[el camino al caos]\label{prob:in-mapa}
\etiquetaM{} (a) Para el mapa \eqref{eq:in-mapa} muestre que $F'(P^*)=2-R$ en el punto fijo no trivial y concluya que pierde la estabilidad en $R=3$. (b) Con \cmd{C1\_Ej2\_Logistica}, encuentre aproximadamente el valor de $R$ en el que aparece el ciclo de período 4. (c) ¿Qué pasa con las iteraciones si $R>4$ y $P[0]\in(0,K)$?
\end{problema}

\begin{problema}[la ley lineal que crea equilibrios]\label{prob:in-cubico}
\etiquetaM{} Para $\dot x=x+u^3$ con $u=-kx$: (a) determine la región de atracción de cada uno de los tres equilibrios del lazo cerrado. (b) ¿Hay algún $k$ que estabilice el origen? (c) ¿Es $u=-(kx)^{1/3}$ una función Lipschitz? ¿Importa, si el lazo cerrado es lineal?
\end{problema}

\begin{problema}[lo que el integrador no corrige]\label{prob:in-pid}
En el \cref{ej:in-pid}: (a) muestre que $y^*=r-w$ para todo $k_I>0$ con que el lazo sea estable. (b) Si $v$ fuera una rampa, $v=at$, ¿seguiría siendo cero el error debido a $v$? (c) ¿Qué habría que cambiar para que $y\to r$?
\end{problema}

\begin{problema}[el período del péndulo]\label{prob:in-periodo}
\etiquetaEstrella{} (a) Con la conservación de la energía, muestre que el período del péndulo sin fricción de amplitud $\theta_0$ es $T=\frac{4}{\omega_0}\int_0^{\pi/2}\bigl(1-k^2\sin^2\varphi\bigr)^{-1/2}\dif\varphi$ con $k=\sin(\theta_0/2)$. (Pista: el cambio $\sin(\theta/2)=k\sin\varphi$.) (b) \etiquetaM{} Verifique la \cref{tab:in-periodo}. (c) ¿Qué pasa con $T$ cuando $\theta_0\to\pi$?
\end{problema}

%% ---------------------------------------------------------------------
\section{Soluciones seleccionadas}
%% ---------------------------------------------------------------------

\paragraph{Problema~\ref{prob:in-equilibrios}}
(a) $0$, $\pm1$: tres equilibrios aislados, imposible en un sistema lineal. (b) Solo $0$, y las soluciones con $x_0>0$ escapan en tiempo finito. (c) $k\pi$, infinitos. (d) $0$ y $r$; en $r=0$ chocan e intercambian la estabilidad (una bifurcación transcrítica, \cref{cap:bifurcaciones}). Ninguno puede tener ciclos límite ni caos: son escalares (\cref{cap:unidimensionales}).

\paragraph{Problema~\ref{prob:in-duffing}}
$x_1=x$, $x_2=\dot x$: $\dot x_1=x_2$, $\dot x_2=-\delta x_2-\alpha x_1-\beta x_1^3+\gamma\cos\omega t$; variante en el tiempo si $\gamma\neq0$. Con $\gamma=0$ los equilibrios son $(x_1,0)$ con $x_1(\alpha+\beta x_1^2)=0$: uno si $\alpha\beta>0$ (o $\beta=0$, $\alpha\neq0$) y tres, $0$ y $\pm\sqrt{-\alpha/\beta}$, si $\alpha\beta<0$. La trampa: con $\gamma\neq0$ la pregunta ``¿cuántos equilibrios?'' no tiene sentido, porque $f(t,x)=0$ para todo $t$ no tiene solución.

\paragraph{Problema~\ref{prob:in-reescalar}}
(a) $\dif z/\dif\tau=\frac1{RK}\dot P=z(1-z)$. (b) $\dif z/\dif\tau=z(1-z)-h$ con $h=H/(RK)$: un solo parámetro, la cosecha medida en unidades de $RK$. En el \cref{cap:unidimensionales} se ve que la población se extingue si $h>1/4$, es decir, si la cosecha supera el crecimiento máximo de la población, $H>RK/4$.

\paragraph{Problema~\ref{prob:in-brockett}}
(a) $f(x,u)=(x_2,u)$: dado $y=(y_1,y_2)$ basta $x_2=y_1$, $u=y_2$, pequeños si $y$ lo es. La condición se cumple, pero es solo \emph{necesaria}: de (a) no se concluye que exista la ley (aunque existe: el PD). (b) Falla, como en la \cref{sec:in-brockett}. (c) El integrador de Brockett no es estabilizable con $u=k(x)$ suave.

\paragraph{Problema~\ref{prob:in-mapa}}
(a) $F'(P)=R(1-2P/K)$ y $P^*=K(1-1/R)$ dan $F'(P^*)=R(2/R-1)=2-R$; $\abs{2-R}<1$ si $1<R<3$. (b) Cerca de $R=1+\sqrt6\approx3.449$. (c) El máximo de $F$ es $RK/4>K$: algunas iteraciones salen de $[0,K]$, la población se vuelve negativa y diverge. El modelo deja de tener sentido físico.

\paragraph{Problema~\ref{prob:in-cubico}}
(a) $(0,\infty)$ va a $k^{-3/2}$ y $(-\infty,0)$ a $-k^{-3/2}$; el origen solo se atrae a sí mismo. (b) No: $f'(0)=1$ para todo $k$. (c) No es Lipschitz en $0$ (pendiente infinita), pero el lazo cerrado $\dot x=(1-k)x$ sí lo es, y es lo que importa para la existencia y la unicidad. En la práctica importa otra cosa: cerca de $0$ la ley pide ganancias enormes, y un ruido de medición pequeño produce entradas grandes.

\respuestaspausas

